\begin{apartat}{1}
Escriviu l'equació del balanç energètic en l'efecte fotoelèctric. Digueu el significat de cada un dels termes i deduïu l'expressió de la longitud d'ona llindar a partir, únicament, de la funció de treball (treball d'extracció) del metall i de constants universals.
\end{apartat}

\begin{apartat}{1}
Il·luminem una placa de sodi (funció de treball, $2,36\ \mathrm{eV}$) amb radiació de freqüència de $660\ \mathrm{THz}$. Calculeu l'energia cinètica màxima dels fotoelectrons que s'emetran i el potencial de frenada necessari per a aturar-los.
\end{apartat}

\dades{%
  Constant de Planck, $h = 6,63\times 10^{-34}\ \mathrm{J\,s}$.\\
  Velocitat de la llum, $c = 3,00\times 10^{8}\ \mathrm{m\,s^{-1}}$.\\
  $1\ \mathrm{eV} = 1,60\times 10^{-19}\ \mathrm{J}$.\\
  Càrrega de l'electró, $q_\mathrm{e} = -1,60\times 10^{-19}\ \mathrm{C}$.}
